\section{Support Is Not Identity}
A component can be necessary for maintaining a system without being identical to the system-level property.

\textbf{Proposition 1 (Support is not identity).} If component $x$ is causally necessary for maintaining $B$ under a specified condition, it does not follow that $x\equiv B$ or $x\equiv C_B$.

A heart can be necessary for circulation without being the organism. Particular neural structures can be necessary for wakefulness, breathing, memory, or reportability without being identical to the continuing individual. Organ replacement makes the distinction visible: if a component is replaced while organism-level continuity is preserved, the individual has not been recreated merely because the implementation changed.

Temporary functional collapse is likewise not identical to terminal discontinuity. Let $\mathcal I$ be a declared family of admissible rescue interventions. Define the terminal boundary
\begin{equation}
\tau_B^*=\inf\left\{t:\substack{\text{no admissible trajectory under }\mathcal I\\
\text{preserves }B\text{ as the same viable individual}}\right\}.
\end{equation}
The criterion is explicitly counterfactual and intervention-relative. Successful resuscitation is evidence that the terminal boundary had not been crossed under the declared intervention set; it is not a proof that all imaginable identity questions are settled.

\section{Biological Individuality Is a Dependency, Not a Hidden Escape}
SCTC inherits rather than solves the biological-individuality problem. Fulda treats strong agential autonomy as a criterion of biological individuality \cite{Fulda2023}; Bechtel and Bich analyze physiological regulation from the perspective of maintaining the organism through time \cite{BechtelBich2024}. These frameworks motivate, but do not completely settle, every boundary case.

Accordingly, where biology can identify an integrated, relatively autonomous individual with acceptable intersubjective reliability, SCTC inherits that classification. Where individuality is genuinely indeterminate, as may occur in some cases involving fission, fusion, pregnancy, colonial organization, cryptobiosis, or intensive machine support, SCTC inherits the indeterminacy. It does not replace an unresolved biological boundary with a neurological stipulation.

This matters for brain death. The 2023 AAN/AAP/CNS/SCCM guideline provides clinical criteria for death by neurologic criteria \cite{Greer2023}. Shewmon documented prolonged somatic survival after formal brain-death diagnoses and argued that absence of brain function does not by itself establish loss of all somatic integrative unity \cite{Shewmon1998}. SCTC is not medical or legal guidance. Its theoretical commitment is narrower:

\textbf{Consequence 1.} If a permanently brain-nonfunctional body still satisfies the best operational criteria for one viable integrated biological individual, then $L_B=1$ and therefore $C_B=1$ under SCTC. If the remaining physiology no longer constitutes one biological individual, then $L_B=0$ and $C_B=0$. If biological individuality is indeterminate, SCTC is indeterminate.

\section{Rival Assignments Must Be Named, Not Smuggled In}
SCTC is not the only coherent assignment.

\subsection{Experience-centered rival}
The nearest naming rival is
\begin{equation}
\mathcal R_E:\qquad C:=E,
\end{equation}
with responsiveness, connectedness, memory, access, reportability, and other variables decomposed. Sanders et al. instantiate this style explicitly \cite{Sanders2012}. Shared success on anesthesia or dissociation does not adjudicate SCTC, because both assignments predict those separations.

\subsection{Capacity-transition rivals}
A biologically substantive rival places the transition to minimal consciousness at a specific capacity rather than at organismal viability. Unlimited Associative Learning (UAL), for example, has been proposed as a transition marker for minimal consciousness or sentience \cite{Bronfman2016,Birch2020}. SCTC does not deny that such capacities may mark the emergence of sentience or phenomenal experience. It denies that a capacity marker is thereby the persistence condition of the biological individual.

The disagreement is therefore explicit:
\begin{align}
\mathcal R_\Gamma &: C:=L_B,\\
\mathcal R_E &: C:=E,\\
\mathcal R_T &: C:=T \quad\text{for a proposed capacity transition }T.
\end{align}
No rival may be used as an unstated definition while simultaneously being offered as an objection.

